\documentclass[11pt,a4paper]{article}
\usepackage[a4paper,top=1.8cm,bottom=1.6cm,left=2.3cm,right=2.3cm]{geometry}
\usepackage{lmodern}
\usepackage[T1]{fontenc}
\usepackage{enumitem}
\usepackage[hidelinks]{hyperref}

\linespread{1.0}
\pagestyle{empty}
\setlength{\parindent}{0pt}
\setlength{\parskip}{6pt}
\emergencystretch=1em

\begin{document}

\begin{flushright}
Dr. Vipin Gupta\\
Department of Mathematics\\
Gurugram University, Gurugram 122505, India\\
\href{mailto:vipin@gurugramuniversity.ac.in}{vipin@gurugramuniversity.ac.in}\\[6pt]
\today
\end{flushright}

To\\
The Editor-in-Chief\\
Applied Mathematics and Mechanics (English Edition)

{\raggedright
\textbf{Subject:} Submission of the manuscript ``Plane-Strain Thermal Stress and
Piezoelectric-Pyroelectric Coupling in a Tetragonal BaTiO$_3$ Microcantilever''
\par}

Dear Editor,

We are pleased to submit this manuscript for consideration as a research paper in
\emph{Applied Mathematics and Mechanics (English Edition)}. The work examines the
open-circuit voltage of a tetragonal BaTiO$_3$ microcantilever subjected to a steady
temperature difference through its thickness, with emphasis on the interaction between
anisotropic thermal expansion, plane-strain compatibility, piezoelectricity, and
pyroelectricity.

The manuscript's main contributions are:
\begin{itemize}[leftmargin=1.5em,itemsep=2pt,topsep=2pt]
\item a plane-strain thermal-stress reduction derived from the full three-dimensional
thermal-expansion tensor, showing that the simplified reduction underestimates the
axial and through-thickness thermal stresses by 35.05\% and 57.75\%, respectively;
\item a finite-element formulation based on a single electric enthalpy that treats the
elastic, piezoelectric, dielectric, thermal, and primary pyroelectric contributions
consistently, using a Bogner--Fox--Schmit element for the mechanical field;
\item the finding that the thermal-stress reduction changes the predicted
piezoelectric-thermal voltage polarity, from $-0.22$~V to $+0.19$~V, together with
mesh-converged open-circuit voltages of $-3.43865$~V for the pyroelectric-only case,
$-1.49960$~V for the coupled case, and $+0.19321$~V for the piezoelectric-thermal
case; and
\item verification against the exact one-dimensional pyroelectric limit, mesh
convergence, and algebraic and variational consistency checks.
\end{itemize}

The results highlight how anisotropic material data and mechanical constraints can
change not only the magnitude but also the polarity of a coupled-field response. We
believe that the manuscript is relevant to the journal's interests in applied
mathematics, mechanics, and computational methods for coupled physical systems.

The manuscript is original, has not been published previously, and is not under
consideration for publication elsewhere. No potential conflict of interest was
reported by the authors.

Thank you for considering our submission.

Sincerely,\\[2pt]
Dr. Vipin Gupta\\
Corresponding author, on behalf of all authors

\end{document}
